\documentclass{article}
\usepackage[utf8]{inputenc}
\usepackage{a4wide}

\begin{document}

\section*{Bridge  {CT}}

ในเมืองแห่งหนึ่ง มีแม่น้ำใหญ่ผ่านกลางเมืองทำให้เมืองนี้แบ่งอำเภอออกเป็น\texttt{2}ฝั่ง คือ

ฝั่งเหนือและฝั่งใต้

ในเมืองแห่งนี้มีทั้งสิ้น\texttt{n}ธนาคาร\textbf{โดยทุกธนาคารจะมีสาขาหนึ่งอยู่ในฝั่งเหนือและมีอีกสาขาหนึ่งอยู่ในเมืองฝั่งใต้เสมอ}(ธนาคารทั้งในฝั่งเหนือและฝั่งใต้จะเรียงตามแนวแกน\texttt{x}โดยเรียงตามตำแหน่งในแนวแกน\texttt{x}ของธนาคารแต่ละสาขา)\textbf{รัฐบาลต้องการสร้างสะพานเพื่อเชื่อมธนาคารเดียวกันระหว่างฝั่งเหนือกับฝั่งใต้ให้ได้จำนวนมากที่สุด}โดยไม่ให้เกิดการตัดกันของสะพาน

(คือไม่ต้องการให้มีการสร้างสะพานระหว่างสะพาน)\\



\hfill\break



\textbf{ตัวอย่าง}



\textbf{input:}



\begin{quote}

\textbf{5}



\textbf{3}



\textbf{1}



\textbf{5}



\textbf{2}



\textbf{4}



\textbf{1}



\textbf{2}



\textbf{5}



\textbf{4}



\textbf{3}

\end{quote}



\textbf{output:}



\begin{quote}

\textbf{3}

\end{quote}



\hfill\break



\includegraphics[width=3.08333in,height=2.08333in,alt={image.png}]{/public/upload/6173779b99.png}\\



{\textbf{อธิบายตัวอย่าง}}{\hfill\break

}ในตัวอย่างนี้ มีธนาคารทั้งหมด 5 ธนาคาร ในฝั่งเหนือธนาคารจะเรียงตามลำดับในแกน x คือ

3, 1, 5, 2, 4 และในฝั่งใต้มีลำดับของธนาคารคือ 1, 2, 5, 4, 3

ซึ่งสามารถสร้างสะพานได้มากที่สุด 3 สะพาน โดยคำตอบที่เป็นไปได้คือสะพานเชื่อมธนาคารที่ 1,

สะพานเชื่อมธนาคารที่ 2, และสะพานเชื่อมธนาคารที่ 4 หรืออาจจะสร้างสะพานที่เชื่อมธนาคารที่

1 แทนสะพานเชื่อมธนาคารที่ 2 ก็ได้ แต่ไม่สามารถสร้างทั้งสองสะพานเชื่อมธนาคารที่ 2

และสะพานเชื่อมธนาคารที่ 5 พร้อมกันได้ เนื่องจากสะพานจะตัดกัน\\



\subsection*{Input}
\ul{บรรทัดแรก}รับจำนวนเต็ม\texttt{N}\texttt{(1\ \textless{}=\ N\ \textless{}=\ 3000)}แทนจำนวนธนาคารทั้งหมดในเมือง\\



{\ul{บรรทัดที่ 2 ถึง

N+1}}รับหมายเลขของธนาคารในฝั่งเหนือโดยเรียงตามพิกัดในแนวแกน\texttt{x}โดยบรรทัดที่\texttt{1\ +\ I}จะรับข้อมูลจำนวนเต็มหนึ่งจำนวนแทนหมายเลขของธนาคารที่อยู่ในฝั่งเหนือและอยู่ในอันดับที่\texttt{I}เมื่อเรียงตามตำแหน่งในแนวแกน\texttt{x}(ข้อมูลทั้งหมด\texttt{n}จำนวนนี้จะไม่ซ้ำกันและมีค่าตั้งแต่\texttt{1}ถึง\texttt{N})\\



\ul{บรรทัดที่ N + 2 ถึง 2N +

1}รับหมายเลขของธนาคารในฝั่งใต้โดยเรียงตามพิกัดในแนวแกน\texttt{x}โดยบรรทัดที่\texttt{N\ +\ 1\ +\ I}จะรับข้อมูลจำนวนเต็มหนึ่งจำนวนแทนหมายเลขของธนาคารที่อยู่ในฝั่งใต้และอยู่ในอันดับที่\texttt{I}เมื่อเรียงตามตำแหน่งในแนวแกน\texttt{x}(ข้อมูลทั้ง\texttt{N}จำนวนนี้จะไม่ซ้ำกันและมีค่าตั้งแต่\texttt{1}ถึง\texttt{N})\\



\subsection*{Output}
ให้แสดงตัวเลขจำนวนเต็มหนึ่งจำนวน

ที่แทนจำนวนสะพานที่มากที่สุดที่สามารถสร้างเชื่อมธนาคารทั้งสองฝั่งได้ โดยไม่มีสะพานที่ตัดกัน\\



\end{document}
